\heading{理论力学-26春-期末}
\noteinfo{题目来源于考后回忆，答案由AI生成。任课教师：邵立晶。}
\newcommand{\ii}{\mathrm{i}}
\newcommand{\ee}{\mathrm{e}}
\begin{question}[含辅助变量的拉格朗日量]
考察包含辅助变量 $\chi$ 的拉格朗日量
\[
L(q,\dot q,\chi)=\frac12 m\dot q^{2}-\frac12 k q^{2}+a\chi\dot q+b\chi^{2}.
\]

\begin{enumerate}[label=(\roman*)]
  \item 写出辅助变量 $\chi$ 满足的 Euler--Lagrange 方程；
  \item 求解 $\chi$，并代入得到等效拉格朗日量 $L_{\text{eff}}(q,\dot q)$；
  \item 分析当参数满足何条件时，会出现``负质量效应''，导致运动不稳定（即``快子不稳定''）。
\end{enumerate}

\begin{solution}

\textbf{(i)} $\chi$ 的 Euler--Lagrange 方程为
\[
\frac{\partial L}{\partial\chi}-\frac{\dd}{\dd t}\frac{\partial L}{\partial\dot\chi}=0.
\]
计算偏导数：
\[
\frac{\partial L}{\partial\chi}=a\dot q+2b\chi,\qquad
\frac{\partial L}{\partial\dot\chi}=0,
\]
故 $\chi$ 满足代数方程（非动力学变量）
\[
a\dot q+2b\chi=0.
\]

\textbf{(ii)} 由上式立得
\[
\chi=-\frac{a}{2b}\,\dot q.
\]
代入原拉格朗日量：
\begin{align*}
L_{\text{eff}}
&=\frac12 m\dot q^{2}-\frac12 k q^{2}
+a\Bigl(-\frac{a}{2b}\dot q\Bigr)\dot q
+b\Bigl(-\frac{a}{2b}\dot q\Bigr)^{\!2} \\[2pt]
&=\frac12 m\dot q^{2}-\frac12 k q^{2}
-\frac{a^{2}}{2b}\dot q^{2}+\frac{a^{2}}{4b}\dot q^{2} \\[2pt]
&=\frac12\Bigl(m-\frac{a^{2}}{2b}\Bigr)\dot q^{2}
-\frac12 k q^{2}.
\end{align*}

\textbf{(iii)} 等效动能项系数为 $m_{\text{eff}}=m-\dfrac{a^{2}}{2b}$。
当 $m_{\text{eff}}<0$ 时出现``负质量效应''，即

\ansmath{
\frac{a^{2}}{2b}>m\quad\Longrightarrow\quad a^{2}>2bm.
}

此时等效动能项变号，体系能量无下界，导致``快子不稳定''。
物理上这对应着辅助场 $\chi$ 与 $q$ 的耦合过强，使得有效惯量为负。
\end{solution}
\end{question}


\begin{question}[Hamilton 矢量场与生成函数]
写出定义在相空间上的函数 $G(q,p)=pq$ 对应的 Hamilton 矢量场，
写出其对应的无穷小正则变换，以及其对应的第二类生成函数。

\begin{solution}

\textbf{Hamilton 矢量场：}
\[
X_{G}
=\frac{\partial G}{\partial p}\frac{\partial}{\partial q}
-\frac{\partial G}{\partial q}\frac{\partial}{\partial p}
=q\,\frac{\partial}{\partial q}-p\,\frac{\partial}{\partial p}.
\]

\textbf{无穷小正则变换：}
设无穷小参数为 $\varepsilon$，则
\[
\delta q=\varepsilon\{q,G\}=\varepsilon\frac{\partial G}{\partial p}
=\varepsilon\,q,\qquad
\delta p=\varepsilon\{p,G\}=-\varepsilon\frac{\partial G}{\partial q}
=-\varepsilon\,p.
\]
即 $(q,p)\to(q+\varepsilon q,\;p-\varepsilon p)$，这是相空间的标度变换（伸缩变换）。
有限形式为
\[
Q=\ee^{\varepsilon}q,\qquad P=\ee^{-\varepsilon}p.
\]

\textbf{第二类生成函数：}
对无穷小正则变换，第二类生成函数的形式为
\[
F_{2}(q,P)=qP+\varepsilon\,G(q,P).
\]
代入 $G(q,P)=qP$ 得
\[
F_{2}(q,P)=qP+\varepsilon\,qP=(1+\varepsilon)qP.
\]
相应地，有限形式的生成函数为

\ansmath{
F_{2}(q,P)=\ee^{\varepsilon}qP,
}
其中 $\varepsilon$ 为变换参数。验证：
\[
p=\frac{\partial F_{2}}{\partial q}=\ee^{\varepsilon}P,\qquad
Q=\frac{\partial F_{2}}{\partial P}=\ee^{\varepsilon}q.
\]

\end{solution}
\end{question}


\begin{question}[阻尼振子]
考察阻尼振子满足的拉格朗日量
\[
L(q,\dot q,t)=\ee^{2\lambda t}
\Bigl(\frac12 m\dot q^{2}-\frac12 m\omega^{2}q^{2}\Bigr).
\]

\begin{enumerate}[label=(\roman*)]
  \item 写出体系的哈密顿量 $H(q,p,t)$；
  \item 计算在第二类生成函数 $F_{2}(q,P,t)=\ee^{\lambda t}qP$ 对应的正则变换下得到的哈密顿量 $K(Q,P,t)$；
  \item 取 $\lambda=\omega$，使用 Hamilton--Jacobi 方程求解 $q(t)$ 和 $p(t)$。
\end{enumerate}

\begin{solution}

\textbf{(i)} 正则动量：
\[
p=\frac{\partial L}{\partial\dot q}=m\ee^{2\lambda t}\dot q
\quad\Longrightarrow\quad
\dot q=\frac{p}{m}\,\ee^{-2\lambda t}.
\]

哈密顿量：
\begin{align*}
H(q,p,t)&=p\dot q-L
=p\cdot\frac{p}{m}\ee^{-2\lambda t}
-\ee^{2\lambda t}\Bigl(\frac12 m\frac{p^{2}}{m^{2}}\ee^{-4\lambda t}
-\frac12 m\omega^{2}q^{2}\Bigr) \\[2pt]
&=\frac{p^{2}}{m}\ee^{-2\lambda t}
-\frac{p^{2}}{2m}\ee^{-2\lambda t}
+\frac12 m\omega^{2}\ee^{2\lambda t}q^{2} \\[2pt]
&=\frac{p^{2}}{2m}\,\ee^{-2\lambda t}
+\frac12 m\omega^{2}\ee^{2\lambda t}q^{2}.
\end{align*}

\textbf{(ii)} 由 $F_{2}=\ee^{\lambda t}qP$：
\[
p=\frac{\partial F_{2}}{\partial q}=\ee^{\lambda t}P,\qquad
Q=\frac{\partial F_{2}}{\partial P}=\ee^{\lambda t}q.
\]
因此
\[
q=\ee^{-\lambda t}Q,\qquad p=\ee^{\lambda t}P.
\]

新哈密顿量：
\[
K(Q,P,t)=H(q,p,t)+\frac{\partial F_{2}}{\partial t}.
\]
计算 $\partial F_{2}/\partial t=\lambda\ee^{\lambda t}qP=\lambda QP$。

将 $q,p$ 用 $Q,P$ 表达代入 $H$：
\begin{align*}
H&=\frac{(\ee^{\lambda t}P)^{2}}{2m}\,\ee^{-2\lambda t}
+\frac12 m\omega^{2}\ee^{2\lambda t}(\ee^{-\lambda t}Q)^{2} \\[2pt]
&=\frac{P^{2}}{2m}+\frac12 m\omega^{2}Q^{2}.
\end{align*}
故

\ansmath{
K(Q,P,t)=\frac{P^{2}}{2m}+\frac12 m\omega^{2}Q^{2}+\lambda QP.
}

\textbf{(iii)} 取 $\lambda=\omega$，则新哈密顿量为
\[
K(Q,P)=\frac{P^{2}}{2m}+\frac12 m\omega^{2}Q^{2}+\omega QP.
\]

$K$ 不含时（时间显无关），故可令主函数 $S(Q,\alpha,t)=W(Q,\alpha)-\alpha t$，
其中 $\alpha$ 为新的守恒动量。Hamilton--Jacobi 方程为
\[
\frac{\partial S}{\partial t}+K\!\Bigl(Q,\frac{\partial S}{\partial Q}\Bigr)=0
\quad\Longrightarrow\quad
\frac{1}{2m}\Bigl(\frac{\partial W}{\partial Q}\Bigr)^{\!2}
+\frac12 m\omega^{2}Q^{2}+\omega Q\frac{\partial W}{\partial Q}
=\alpha.
\]

这是 $\partial W/\partial Q$ 的二次方程：
\[
\Bigl(\frac{\partial W}{\partial Q}\Bigr)^{\!2}
+2m\omega Q\frac{\partial W}{\partial Q}
+m^{2}\omega^{2}Q^{2}-2m\alpha=0.
\]

配方得
\[
\Bigl(\frac{\partial W}{\partial Q}+m\omega Q\Bigr)^{\!2}=2m\alpha
\quad\Longrightarrow\quad
\frac{\partial W}{\partial Q}=-m\omega Q\pm\sqrt{2m\alpha}.
\]

积分（取 $+$ 号，$-$ 号等价）得
\[
W(Q,\alpha)=-\frac12 m\omega Q^{2}+\sqrt{2m\alpha}\,Q.
\]

故主函数为
\[
S(Q,\alpha,t)=-\frac12 m\omega Q^{2}+\sqrt{2m\alpha}\,Q-\alpha t.
\]

新的守恒坐标 $\beta$ 为
\[
\beta=\frac{\partial S}{\partial\alpha}
=\frac{mQ}{\sqrt{2m\alpha}}-t
=Q\sqrt{\frac{m}{2\alpha}}-t.
\]

因此
\[
Q=\sqrt{\frac{2\alpha}{m}}\,(t+\beta).
\]

新哈密顿量为零（$K'=0$），$\alpha,\beta$ 均为常数。由 $P=\partial S/\partial Q$ 得
\[
P=\frac{\partial S}{\partial Q}=-m\omega Q+\sqrt{2m\alpha}
=\sqrt{2m\alpha}\bigl[1-\omega(t+\beta)\bigr].
\]

回到原变量 ($q=\ee^{-\omega t}Q$，$p=\ee^{\omega t}P$)：

\ansmath{
q(t)=\ee^{-\omega t}\sqrt{\frac{2\alpha}{m}}\,(t+\beta),\qquad
p(t)=\sqrt{2m\alpha}\,\ee^{\omega t}\bigl[1-\omega(t+\beta)\bigr].
}

引入 $A=\sqrt{2\alpha/m}$，$B=\beta$，则简写为

\ansmath{
q(t)=\ee^{-\omega t}A\,(t+B),\qquad
p(t)=m\ee^{\omega t}A\bigl[1-\omega(t+B)\bigr].
}

这正是$\lambda=\omega$（临界阻尼）情形下的精确解。与原轨迹方程
$\ddot Q=0\to Q=C_{1}t+C_{2}$ 所得结果完全一致。

\end{solution}
\end{question}


\begin{question}[隆格楞次矢量]
隆格楞次矢量定义为
\[
\vec A=\vec r\times\vec J-\alpha m\hat r,
\]
其中 $\vec J=\vec r\times\vec p$ 为角动量，$\hat r=\vec r/r$ 为径向单位矢量，$\alpha$ 为常参数。
计算 $\{A^{i},J^{j}\}$（Poisson 括号）。

\begin{solution}

在直角坐标下，$J^{j}=\varepsilon^{jmn}r_{m}p_{n}$。
首先将 $A^{i}$ 用坐标和动量表出：
\begin{align*}
A^{i}&=\varepsilon^{ikl}r_{k}J^{l}-\alpha m\frac{r^{i}}{r} \\
&=\varepsilon^{ikl}\varepsilon^{lmn}r_{k}r_{m}p_{n}
-\alpha m\frac{r^{i}}{r}.
\end{align*}

利用恒等式 $\varepsilon^{ikl}\varepsilon^{lmn}
=\varepsilon^{lik}\varepsilon^{lmn}
=\delta^{im}\delta^{kn}-\delta^{in}\delta^{km}$：
\[
\varepsilon^{ikl}\varepsilon^{lmn}r_{k}r_{m}p_{n}
=(\delta^{im}\delta^{kn}-\delta^{in}\delta^{km})r_{k}r_{m}p_{n}
=r_{n}r^{i}p_{n}-r_{m}r_{m}p^{i}
=r^{i}(\vec r\cdot\vec p)-r^{2}p^{i}.
\]

因此
\[
A^{i}=r^{i}(\vec r\cdot\vec p)-r^{2}p^{i}-\alpha m\frac{r^{i}}{r}.
\]

现在计算 Poisson 括号 $\{A^{i},J^{j}\}$。利用 $J^{j}$ 是旋转生成元这一事实，
$\{A^{i},J^{j}\}$ 应当体现矢量 $\vec A$ 的旋转性质：

\ansmath{
\{A^{i},J^{j}\}=\varepsilon^{ijk}A^{k}.
}

这一结果可直接验证。以 $j=1$ 为例（$J^{1}=yp_{z}-zp_{y}$）：

计算 $\{A^{2},J^{1}\}\stackrel{?}{=}\varepsilon^{213}A^{3}=-A^{3}$。
由 $\{r^{i},J^{j}\}=\varepsilon^{ijk}r^{k}$，$\{p^{i},J^{j}\}=\varepsilon^{ijk}p^{k}$，
以及 $A^{i}$ 是 $r^{i},p^{i}$ 的矢量组合（$A^{i}$ 中各项均是矢量的分量），
根据 Poisson 括号的链式法则可得上述结果。

\textbf{验证要点：}
\begin{itemize}
  \item $\{r^{i}(\vec r\cdot\vec p),J^{j}\}
  =\{r^{i},J^{j}\}(\vec r\cdot\vec p)+r^{i}\{\vec r\cdot\vec p,J^{j}\}
  =\varepsilon^{ijk}r^{k}(\vec r\cdot\vec p)+r^{i}\cdot0$（因 $\vec r\cdot\vec p$ 是转动标量）；
  \item $\{r^{2}p^{i},J^{j}\}
  =\{r^{2},J^{j}\}p^{i}+r^{2}\{p^{i},J^{j}\}
  =0\cdot p^{i}+r^{2}\varepsilon^{ijk}p^{k}
  =\varepsilon^{ijk}r^{2}p^{k}$；
  \item $\bigl\{\dfrac{r^{i}}{r},J^{j}\bigr\}
  =\dfrac{1}{r}\{r^{i},J^{j}\}+r^{i}\bigl\{\dfrac{1}{r},J^{j}\bigr\}
  =\dfrac{1}{r}\varepsilon^{ijk}r^{k}+0$。
\end{itemize}

将各部分组合即得
\[
\{A^{i},J^{j}\}
=\varepsilon^{ijk}\bigl[r^{k}(\vec r\cdot\vec p)-r^{2}p^{k}-\alpha m\tfrac{r^{k}}{r}\bigr]
=\varepsilon^{ijk}A^{k}.
\]

这正表明 $\vec A$ 在空间旋转下如同一个三维矢量变换，
它是 $\vec r$ 和 $\vec p$ 构成的唯一与 $\vec J$ 对合的矢量守恒量（在 Kepler 问题中）。
\end{solution}
\end{question}


\begin{question}[场论——哈密顿原理与正则方程]

\begin{enumerate}[label=(\roman*)]
  \item 一个 $3+1$ 维体系 $(t,x,y,z)$ 的拉氏量密度为
  $\mathcal{L}(\eta,\eta_{,\mu},\eta_{,\mu\nu},x^{\mu})$，
  使用哈密顿原理给出 $\eta$ 满足的场方程；
  \item 一个 $1+1$ 维体系 $(t,x)$ 的哈密顿量密度为
  $\mathcal{H}(\eta,\eta_{,x},\pi,\pi_{,x},\pi_{,xx})$，
  给出其对应的正则方程组。
\end{enumerate}

\begin{solution}

\textbf{(i)} 作用量为 $S[\eta]=\displaystyle\int\dd^{4}x\;\mathcal{L}
(\eta,\eta_{,\mu},\eta_{,\mu\nu},x^{\mu})$。
哈密顿原理要求 $\delta S=0$：
\begin{align*}
\delta S&=\int\dd^{4}x\;\Bigl[
\frac{\partial\mathcal{L}}{\partial\eta}\delta\eta
+\frac{\partial\mathcal{L}}{\partial\eta_{,\mu}}\delta\eta_{,\mu}
+\frac{\partial\mathcal{L}}{\partial\eta_{,\mu\nu}}\delta\eta_{,\mu\nu}
\Bigr] \\[2pt]
&=\int\dd^{4}x\;
\frac{\partial\mathcal{L}}{\partial\eta}\delta\eta
-\int\dd^{4}x\;\partial_{\mu}\!\Bigl(\frac{\partial\mathcal{L}}{\partial\eta_{,\mu}}\Bigr)\delta\eta
+\int\dd^{4}x\;\partial_{\mu}\partial_{\nu}\!\Bigl(
\frac{\partial\mathcal{L}}{\partial\eta_{,\mu\nu}}\Bigr)\delta\eta,
\end{align*}
其中用到了两次分部积分，并设边界上 $\delta\eta$ 及其一阶导数为零。

令 $\delta S=0$ 对任意 $\delta\eta$ 成立，得到

\ansmath{
\frac{\partial\mathcal{L}}{\partial\eta}
-\partial_{\mu}\Bigl(\frac{\partial\mathcal{L}}{\partial\eta_{,\mu}}\Bigr)
+\partial_{\mu}\partial_{\nu}\Bigl(\frac{\partial\mathcal{L}}{\partial\eta_{,\mu\nu}}\Bigr)=0.
}

这是含二阶导数的广义 Euler--Lagrange 方程。

\textbf{(ii)} 正则变量为 $\eta$ 和 $\pi$，哈密顿量为
$\displaystyle H=\int\dd x\;\mathcal{H}(\eta,\eta_{,x},\pi,\pi_{,x},\pi_{,xx})$。

由变分原理 $\delta\int\dd t\dd x\,(\pi\dot\eta-\mathcal{H})=0$，得正则方程：
\[
\dot\eta=\frac{\delta H}{\delta\pi},\qquad
\dot\pi=-\frac{\delta H}{\delta\eta},
\]
其中 $\delta/\delta\pi$ 表示泛函导数。

对 $\mathcal{H}$ 含 $\pi$ 的二阶空间导数情形：

\ansmath{
\begin{aligned}
\dot\eta&=\frac{\partial\mathcal{H}}{\partial\pi}
-\partial_{x}\Bigl(\frac{\partial\mathcal{H}}{\partial\pi_{,x}}\Bigr)
+\partial_{x}^{2}\Bigl(\frac{\partial\mathcal{H}}{\partial\pi_{,xx}}\Bigr),\\[6pt]
\dot\pi&=-\frac{\partial\mathcal{H}}{\partial\eta}
+\partial_{x}\Bigl(\frac{\partial\mathcal{H}}{\partial\eta_{,x}}\Bigr).
\end{aligned}}

泛函导数由 $\displaystyle\frac{\delta H}{\delta\pi}
=\sum_{n=0}^{\infty}(-1)^{n}\partial_{x}^{n}
\frac{\partial\mathcal{H}}{\partial(\partial_{x}^{n}\pi)}$ 给出。

\end{solution}
\end{question}


\begin{question}[复标量电子场]
复标量电子场的拉格朗日量密度为
\[
\mathcal{L}_{0}=\phi^{,\mu}\phi_{,\mu}^{*}-m^{2}\phi\phi^{*},
\]
其对应的 $\mathrm{U}(1)$ Noether 流为
$j^{\mu}=\ii(\phi^{*}\phi^{,\mu}-\phi\phi^{*,\mu})$。
在外场 $A_{\mu}$ 下，拉格朗日量密度会多一项 $j^{\mu}A_{\mu}$，
即 $\mathcal{L}=\mathcal{L}_{0}+j^{\mu}A_{\mu}$。

\begin{enumerate}[label=(\roman*)]
  \item 给出 $\phi,\phi^{*}$ 满足的场方程；
  \item 给出守恒流与对应的守恒定律。
\end{enumerate}

\begin{solution}

\textbf{(i)} 总拉格朗日量密度：
\[
\mathcal{L}
=\partial_{\mu}\phi\,\partial^{\mu}\phi^{*}-m^{2}\phi\phi^{*}
+\ii A_{\mu}\,(\phi^{*}\partial^{\mu}\phi-\phi\,\partial^{\mu}\phi^{*}).
\]

注意第二项中 $A_{\mu}$ 与 $\phi^{*}\partial^{\mu}\phi$ 均为普通的数（经典场），
可交换顺序。

对 $\phi^{*}$ 做变分：
\[
\frac{\partial\mathcal{L}}{\partial\phi^{*}}
=-m^{2}\phi+\ii A_{\mu}\,\partial^{\mu}\phi
=-m^{2}\phi+\ii A^{\mu}\partial_{\mu}\phi,
\]
\[
\frac{\partial\mathcal{L}}{\partial(\partial_{\mu}\phi^{*})}
=\partial^{\mu}\phi-\ii A^{\mu}\phi.
\]

代入 Euler--Lagrange 方程
$\displaystyle\frac{\partial\mathcal{L}}{\partial\phi^{*}}
-\partial_{\mu}\frac{\partial\mathcal{L}}{\partial(\partial_{\mu}\phi^{*})}=0$：
\begin{align*}
0&=-m^{2}\phi+\ii A^{\mu}\partial_{\mu}\phi
-\partial_{\mu}\bigl(\partial^{\mu}\phi-\ii A^{\mu}\phi\bigr) \\
&=-m^{2}\phi+\ii A^{\mu}\partial_{\mu}\phi
-\square\phi+\ii\partial_{\mu}(A^{\mu}\phi) \\
&=-\square\phi-m^{2}\phi
+\ii A^{\mu}\partial_{\mu}\phi
+\ii(\partial_{\mu}A^{\mu})\phi+\ii A^{\mu}\partial_{\mu}\phi \\
&=-\square\phi-m^{2}\phi
+2\ii A^{\mu}\partial_{\mu}\phi+\ii(\partial_{\mu}A^{\mu})\phi.
\end{align*}

整理得 $\phi$ 的场方程

\ansmath{
\square\phi+m^{2}\phi-2\ii A^{\mu}\partial_{\mu}\phi-\ii(\partial_{\mu}A^{\mu})\phi=0.
}

对 $\phi$ 做变分类似可得 $\phi^{*}$ 的共轭场方程：
\[
\square\phi^{*}+m^{2}\phi^{*}+2\ii A^{\mu}\partial_{\mu}\phi^{*}+\ii(\partial_{\mu}A^{\mu})\phi^{*}=0.
\]

\textbf{注：}若希望将方程写成协变导数的紧凑形式，注意到
$(\partial_{\mu}-\ii A_{\mu})^{2}\phi+m^{2}\phi
=\square\phi-2\ii A^{\mu}\partial_{\mu}\phi-\ii(\partial_{\mu}A^{\mu})\phi-A^{2}\phi+m^{2}\phi$。
由于拉氏量只加了 $j^{\mu}A_{\mu}$ 项而未包含 $A^{2}$ 项，
场方程中不含 $A^{2}\phi$，因此不等于通常的 $D_{\mu}D^{\mu}\phi+m^{2}\phi=0$。
这正是``外场''处理的含义——$A_{\mu}$ 是给定的背景场，不参与动力学。

\textbf{(ii)} 总拉格朗日量 $\mathcal{L}=\mathcal{L}_{0}+j^{\mu}A_{\mu}$ 仍具有整体 $\mathrm{U}(1)$ 对称性（$A_{\mu}$ 作为背景场不参与变换）：
\[
\phi\to\ee^{-\ii\alpha}\phi,\qquad
\phi^{*}\to\ee^{+\ii\alpha}\phi^{*},
\]
无穷小形式为 $\delta\phi=-\ii\alpha\phi$，$\delta\phi^{*}=+\ii\alpha\phi^{*}$。

由 Noether 定理，守恒流需对 $\mathcal{L}$ 而非仅 $\mathcal{L}_{0}$ 计算：
\[
J^{\mu}
=\frac{\partial\mathcal{L}}{\partial(\partial_{\mu}\phi)}\,\frac{\delta\phi}{\alpha}
+\frac{\partial\mathcal{L}}{\partial(\partial_{\mu}\phi^{*})}\,\frac{\delta\phi^{*}}{\alpha}.
\]

偏导数为
\[
\frac{\partial\mathcal{L}}{\partial(\partial_{\mu}\phi)}
=\partial^{\mu}\phi^{*}+\ii A^{\mu}\phi^{*},\qquad
\frac{\partial\mathcal{L}}{\partial(\partial_{\mu}\phi^{*})}
=\partial^{\mu}\phi-\ii A^{\mu}\phi.
\]

代入得
\begin{align*}
J^{\mu}
&=\bigl(\partial^{\mu}\phi^{*}+\ii A^{\mu}\phi^{*}\bigr)(-\ii\phi)
+\bigl(\partial^{\mu}\phi-\ii A^{\mu}\phi\bigr)(+\ii\phi^{*}) \\[2pt]
&=-\ii\phi\,\partial^{\mu}\phi^{*}+A^{\mu}|\phi|^{2}
+\ii\phi^{*}\partial^{\mu}\phi+A^{\mu}|\phi|^{2} \\[2pt]
&=\ii\bigl(\phi^{*}\partial^{\mu}\phi-\phi\,\partial^{\mu}\phi^{*}\bigr)
+2A^{\mu}|\phi|^{2}.
\end{align*}

因此含外场的总守恒流为

\ansmath{
J^{\mu}=j^{\mu}+2A^{\mu}|\phi|^{2},
\qquad
j^{\mu}\equiv\ii(\phi^{*}\partial^{\mu}\phi-\phi\,\partial^{\mu}\phi^{*}).
}

\textbf{守恒律验证：}由第 (i) 问得到的场方程
\begin{align*}
\square\phi+m^{2}\phi-2\ii A^{\mu}\partial_{\mu}\phi-\ii(\partial_{\mu}A^{\mu})\phi &=0,\\
\square\phi^{*}+m^{2}\phi^{*}+2\ii A^{\mu}\partial_{\mu}\phi^{*}+\ii(\partial_{\mu}A^{\mu})\phi^{*} &=0,
\end{align*}
计算
\begin{align*}
\partial_{\mu}J^{\mu}&=
\ii\bigl(\partial_{\mu}\phi^{*}\partial^{\mu}\phi+\phi^{*}\square\phi
-\partial_{\mu}\phi\partial^{\mu}\phi^{*}-\phi\square\phi^{*}\bigr)
+2\partial_{\mu}(A^{\mu}|\phi|^{2}) \\[2pt]
&=\ii(\phi^{*}\square\phi-\phi\square\phi^{*})
+2(\partial_{\mu}A^{\mu})|\phi|^{2}+2A^{\mu}\partial_{\mu}|\phi|^{2}.
\end{align*}
代入场方程，并利用 $\partial_{\mu}|\phi|^{2}=\phi^{*}\partial_{\mu}\phi+\phi\,\partial_{\mu}\phi^{*}$，
各项相互抵消，故

\ansmath{
\partial_{\mu}J^{\mu}=0.
}

对应的守恒荷为
\[
Q=\int\dd^{3}x\;J^{0}
=\int\dd^{3}x\;\bigl[\ii(\phi^{*}\dot\phi-\phi\dot\phi^{*})+2A^{0}|\phi|^{2}\bigr],
\qquad \frac{\dd Q}{\dd t}=0.
\]

\textbf{物理意义：}外场 $A_{\mu}$ 存在时，$\mathrm{U}(1)$ 对称性依然保持，
但守恒流 $J^{\mu}$ 相比自由情形多出一项 $2A^{\mu}|\phi|^{2}$，
反映外场对荷密度分布的贡献。

\end{solution}
\end{question}

\clearpage
